\documentclass[11pt,a4paper]{article}

% =========================================================
% PACKAGES
% =========================================================
\usepackage[margin=0.8in]{geometry}
\usepackage{amsmath, amssymb, bm}
\usepackage{booktabs}
\usepackage{enumitem}
\usepackage{hyperref}
\usepackage{xcolor}
\usepackage{array}
\IfFileExists{pdflscape.sty}{\usepackage{pdflscape}}{\usepackage{lscape}}
\usepackage{tabularx}
\usepackage{ragged2e}

\newcolumntype{Y}{>{\RaggedRight\arraybackslash}X}

\hypersetup{
    colorlinks=true,
    linkcolor=blue,
    citecolor=blue,
    urlcolor=blue
}

% =========================================================
% TITLE
% =========================================================
\title{\textbf{Summary Notes on Cycle-Jump Techniques}}
\author{Pruthvi}
\date{\today}

\begin{document}
\maketitle

\tableofcontents
\newpage

% =========================================================
\section{Cojocaru and Karlsson (2006): A Simple Numerical Method of Cycle Jumps for Cyclically Loaded Structures}

\subsection{Full Reference}

Cojocaru, D. and Karlsson, A. M. (2006).
\textit{A simple numerical method of cycle jumps for cyclically loaded structures}.
International Journal of Fatigue, 28, 1677--1689.

\subsection{Main Problem Addressed}

The paper proposes a cycle-jump method for accelerating finite element simulations of cyclically loaded structures whose state evolves slowly from cycle to cycle.

A direct cycle-by-cycle finite element simulation is expensive because each loading cycle may contain many increments and nonlinear iterations. Therefore, simulating all cycles up to failure or long-term evolution may be computationally infeasible.

The main idea is:

\begin{quote}
Compute a few complete finite element cycles, identify the long-term trend of the evolving variables, extrapolate the state over many skipped cycles, and restart the finite element simulation from the extrapolated state.
\end{quote}

\subsection{Basic Cycle-Jump Idea}

The total response of a cyclically loaded structure has two parts:

\begin{itemize}[leftmargin=*]
    \item \textbf{Local cyclic response}: the fast oscillation inside each load cycle.
    \item \textbf{Global evolution}: the slow change of the state from one cycle to the next.
\end{itemize}

The cycle-jump method skips the repeated local cycles and follows the global evolution.

For a state variable $y(t,M)$, where $t$ is time and $M$ is a material point, the method estimates the slow evolution using characteristic values taken at the same point within consecutive cycles.

Examples of possible state variables are:

\begin{itemize}[leftmargin=*]
    \item displacement components,
    \item stress components,
    \item strain components,
    \item equivalent plastic strain,
    \item other state-dependent variables.
\end{itemize}

\subsection{General Algorithm}

The method can be summarized as follows:

\begin{enumerate}[leftmargin=*]
    \item Run a small number of full finite element cycles.
    \item Extract characteristic values of selected state variables at the same relative location in each cycle.
    \item Approximate the global evolution slope of each variable.
    \item Use a control function to determine an admissible jump length.
    \item Extrapolate the variables using the computed jump length.
    \item Impose the extrapolated state as the initial state for the next finite element analysis.
    \item Repeat the procedure.
\end{enumerate}

\subsection{Characteristic Values}

Assume that three consecutive characteristic values of a variable $y$ are available:

\[
P_1(t_1,y_1), \qquad P_2(t_2,y_2), \qquad P_3(t_3,y_3).
\]

Here, $t_1$ is the latest cycle point, while $t_2$ and $t_3$ are from previous cycles.

The cycle length is

\[
\Delta t_{\mathrm{cycle}} = t_1 - t_2 = t_2 - t_3.
\]

The change in the variable over the last two cycles is

\[
\Delta y(t_1) = y_1 - y_2,
\]

\[
\Delta y(t_2) = y_2 - y_3.
\]

The corresponding discrete slopes are

\[
s_{12}(t_1) =
\frac{\Delta y(t_1)}{\Delta t_{\mathrm{cycle}}},
\]

\[
s_{23}(t_2) =
\frac{\Delta y(t_2)}{\Delta t_{\mathrm{cycle}}}.
\]

\subsection{Meaning of the Slopes}

The slope $s_{12}(t_1)$ represents the most recent cycle-to-cycle evolution rate of the variable $y$.

The slope $s_{23}(t_2)$ represents the previous cycle-to-cycle evolution rate.

If

\[
s_{12}(t_1) \approx s_{23}(t_2),
\]

then the global evolution is nearly linear and a larger cycle jump is acceptable.

If

\[
s_{12}(t_1) \not\approx s_{23}(t_2),
\]

then the global evolution is nonlinear and the jump length must be reduced.

\subsection{Control Function for Jump Length}

The predicted slope after a jump is denoted as

\[
s_p\left(t_1+\Delta t^M_{y,\mathrm{jump}}\right).
\]

The control condition is

\[
\frac{
\left|
s_p\left(t_1+\Delta t^M_{y,\mathrm{jump}}\right)
-
s_{12}(t_1)
\right|
}{
\left|s_{12}(t_1)\right|
}
\leq q_y.
\]

Here:

\begin{itemize}[leftmargin=*]
    \item $q_y$ is the user-defined relative error tolerance for variable $y$,
    \item $\Delta t^M_{y,\mathrm{jump}}$ is the locally allowed jump length for variable $y$ at material point $M$,
    \item $s_p$ is the predicted slope after the jump.
\end{itemize}

\subsection{Predicted Slope}

The predicted slope after the jump is obtained by linear extrapolation of the slope:

\[
s_p\left(t_1+\Delta t^M_{y,\mathrm{jump}}\right)
=
s_{12}(t_1)
+
\frac{
s_{12}(t_1)-s_{23}(t_2)
}{
\Delta t_{\mathrm{cycle}}
}
\Delta t^M_{y,\mathrm{jump}}.
\]

This equation predicts how much the evolution rate itself will change after a proposed jump.

\subsection{Allowed Local Jump Length}

Substituting the predicted slope into the control condition gives the allowed local jump length:

\[
\Delta t^M_{y,\mathrm{jump}}
=
q_y \Delta t_{\mathrm{cycle}}
\frac{
\left|s_{12}(t_1)\right|
}{
\left|s_{12}(t_1)-s_{23}(t_2)\right|
}.
\]

This is one of the most important equations in the paper.

\subsubsection*{Interpretation}

If the slopes are almost equal,

\[
\left|s_{12}(t_1)-s_{23}(t_2)\right| \rightarrow 0,
\]

then the denominator becomes small and the allowed jump becomes large.

This means the response is nearly linear and a larger cycle jump is safe.

If the slopes differ strongly, then

\[
\left|s_{12}(t_1)-s_{23}(t_2)\right|
\]

becomes large, and the allowed jump becomes small.

This means the response is nonlinear and the method automatically reduces the cycle jump length.

\subsection{Global Jump Length}

The local jump length is computed for every selected variable and every relevant material point or node.

The global jump length is chosen as the minimum of all local admissible jump lengths:

\[
\Delta t_{\mathrm{jump}}
=
\Delta t_{\mathrm{cycle}}
\left\lfloor
\frac{
\min\left\{\Delta t^M_{y,\mathrm{jump}}\right\}
}{
\Delta t_{\mathrm{cycle}}
}
\right\rfloor.
\]

Here, $\lfloor \cdot \rfloor$ is the floor function.

\subsubsection*{Meaning}

The method uses the most restrictive variable or point to decide the global jump length.

This is conservative but important, because one rapidly evolving point can make a large jump inaccurate.

\subsection{Heun Extrapolation Formula}

After the global jump length is selected, the variable $y$ is extrapolated using a Heun-type integration formula:

\[
y(t_1+\Delta t_{\mathrm{jump}})
=
y(t_1)
+
\frac{1}{2}
\left[
s_{12}(t_1)
+
s_p(t_1+\Delta t_{\mathrm{jump}})
\right]
\Delta t_{\mathrm{jump}}.
\]

Substituting the expression for $s_p$ gives

\[
y(t_1+\Delta t_{\mathrm{jump}})
=
y(t_1)
+
s_{12}(t_1)\Delta t_{\mathrm{jump}}
+
\frac{
\left[
s_{12}(t_1)-s_{23}(t_2)
\right]
(\Delta t_{\mathrm{jump}})^2
}{
2\Delta t_{\mathrm{cycle}}
}.
\]

\subsubsection*{Interpretation}

The first term,

\[
y(t_1),
\]

is the current value.

The second term,

\[
s_{12}(t_1)\Delta t_{\mathrm{jump}},
\]

is the linear extrapolation using the latest slope.

The third term,

\[
\frac{
\left[
s_{12}(t_1)-s_{23}(t_2)
\right]
(\Delta t_{\mathrm{jump}})^2
}{
2\Delta t_{\mathrm{cycle}}
},
\]

corrects the extrapolation using the change in slope between the last two cycles.

Therefore, the method is more accurate than a simple forward Euler extrapolation.

\subsection{Adaptive Number of Full FE Cycles Between Jumps}

After each jump, the method does not immediately jump again. It computes some full finite element cycles before the next jump.

The number of full FE cycles after jump $j$ is denoted by $N_{\mathrm{FEA},j}$.

The empirical update rule is

\[
N_{\mathrm{FEA},j}
=
N_{\mathrm{FEA},j-1}
+
\Delta_+,
\qquad
\text{if }
\frac{\Delta t_{\mathrm{jump},j}}{\Delta t_{\mathrm{cycle}}}
\leq 1,
\]

\[
N_{\mathrm{FEA},j}
=
N_{\mathrm{FEA},j-1}
-
\Delta_-,
\]

if

\[
\frac{\Delta t_{\mathrm{jump},j}}{\Delta t_{\mathrm{cycle}}}
\geq
N_{\mathrm{FEA},j-1}
\]

and

\[
N_{\mathrm{FEA},j-1}-\Delta_- \geq N_{\min}.
\]

Otherwise,

\[
N_{\mathrm{FEA},j}
=
N_{\mathrm{FEA},j-1}.
\]

Here:

\begin{itemize}[leftmargin=*]
    \item $\Delta_+$ increases the number of full FE cycles when jumps become too short,
    \item $\Delta_-$ decreases the number of full FE cycles when jumps are sufficiently long,
    \item $N_{\min}$ is the minimum number of full FE cycles allowed between jumps.
\end{itemize}

\subsection{Why the Control Function Is Necessary}

The paper shows that finite element convergence after a cycle jump does not guarantee that the solution is physically correct.

A large uncontrolled jump may converge numerically but still produce wrong stress and plastic strain fields.

Therefore, the control function is essential because it controls:

\begin{itemize}[leftmargin=*]
    \item the slope change,
    \item the jump length,
    \item the balance between accuracy and computational efficiency.
\end{itemize}

\subsection{Variables Used in the Control Set}

In the examples of the paper, the control set mainly includes displacement components:

\[
u_1, \qquad u_2.
\]

However, the method is general. The control set may include:

\[
\mathcal{C}
=
\{u_i,\ \sigma_{ij},\ \varepsilon_{ij},\ \varepsilon^p_{\mathrm{eq}}, \ldots\}.
\]

The global jump length is then controlled by the most restrictive variable in $\mathcal{C}$.

\subsection{Equivalent Plastic Strain Used for Validation}

The equivalent plastic strain is used as one of the important quantities to assess the accuracy of the cycle-jump result.

For classical plasticity, the equivalent plastic strain is defined as

\[
\varepsilon^p_{\mathrm{eq}}
=
\int d\varepsilon^p_{\mathrm{eq}}
=
\int
\left(
\frac{2}{3}
d\varepsilon^p_{ij}
d\varepsilon^p_{ij}
\right)^{1/2}.
\]

This quantity accumulates with plastic deformation and is useful for checking whether the cycle-jump solution captures progressive inelastic evolution.

\subsection{Computational Efficiency Measure}

The computational efficiency is measured using

\[
R =
\frac{L_{\mathrm{jumps}}}{L_{\mathrm{tot}}},
\]

where:

\begin{itemize}[leftmargin=*]
    \item $L_{\mathrm{jumps}}$ is the total time or number of cycles covered by jumps,
    \item $L_{\mathrm{tot}}$ is the total time or number of cycles considered in the analysis.
\end{itemize}

A larger value of $R$ means that more of the simulation was covered by cycle jumps and fewer cycles were computed directly.

\subsection{Error Measures}

The relative error of a state variable is defined as

\[
\delta E
=
\left|
\frac{
y_{\mathrm{ref}} - y_{\mathrm{jump}}
}{
y_{\mathrm{ref}}
}
\right|
\times 100 \%.
\]

Here:

\begin{itemize}[leftmargin=*]
    \item $y_{\mathrm{ref}}$ is the reference cycle-by-cycle solution,
    \item $y_{\mathrm{jump}}$ is the cycle-jump solution.
\end{itemize}

The average relative error over all integration points is

\[
\overline{\delta E}
=
\frac{1}{N}
\sum_{k=1}^{N}
\delta E_k.
\]

The standard deviation of the relative error is

\[
s
=
\sqrt{
\frac{1}{N}
\sum_{k=1}^{N}
\left(
\delta E_k - \overline{\delta E}
\right)^2
}.
\]

\subsection{Cycle-Jump Procedure in Compact Form}

\begin{enumerate}[leftmargin=*]
    \item Select a control set of variables $\mathcal{C}$.
    \item Run a few full FE cycles.
    \item For each control variable $y$ and material point $M$, extract
    \[
    y_1,\quad y_2,\quad y_3.
    \]
    \item Compute the slopes
    \[
    s_{12}(t_1)
    =
    \frac{y_1-y_2}{\Delta t_{\mathrm{cycle}}},
    \qquad
    s_{23}(t_2)
    =
    \frac{y_2-y_3}{\Delta t_{\mathrm{cycle}}}.
    \]
    \item Compute the locally allowed jump length
    \[
    \Delta t^M_{y,\mathrm{jump}}
    =
    q_y \Delta t_{\mathrm{cycle}}
    \frac{
    |s_{12}(t_1)|
    }{
    |s_{12}(t_1)-s_{23}(t_2)|
    }.
    \]
    \item Compute the global jump length
    \[
    \Delta t_{\mathrm{jump}}
    =
    \Delta t_{\mathrm{cycle}}
    \left\lfloor
    \frac{
    \min\left\{\Delta t^M_{y,\mathrm{jump}}\right\}
    }{
    \Delta t_{\mathrm{cycle}}
    }
    \right\rfloor.
    \]
    \item Extrapolate the variable:
    \[
    y(t_1+\Delta t_{\mathrm{jump}})
    =
    y(t_1)
    +
    s_{12}(t_1)\Delta t_{\mathrm{jump}}
    +
    \frac{
    [s_{12}(t_1)-s_{23}(t_2)]
    (\Delta t_{\mathrm{jump}})^2
    }{
    2\Delta t_{\mathrm{cycle}}
    }.
    \]
    \item Impose the extrapolated state in the FE model.
    \item Run full FE cycles again.
    \item Repeat until the target number of cycles is reached.
\end{enumerate}

\subsection{Implementation Idea in Abaqus}

The paper uses Abaqus with external scripting and user subroutines.

The implementation logic is:

\begin{itemize}[leftmargin=*]
    \item Abaqus/CAE and Python scripting are used to generate and manage the model.
    \item Abaqus solves the complete cycles.
    \item Results are extracted from the output database.
    \item Extrapolated states are written to external files.
    \item User subroutines impose the extrapolated state.
\end{itemize}

The important subroutines mentioned are:

\begin{itemize}[leftmargin=*]
    \item \texttt{SDVINI}: initialization of state variables,
    \item \texttt{DISP}: prescribed extrapolated displacements,
    \item \texttt{UEXPAN}: imposed stress-free growth strain,
    \item \texttt{UMAT}: material behavior and insertion of extrapolated stress/strain information.
\end{itemize}

\subsection{Key Technical Takeaways}

\begin{itemize}[leftmargin=*]
    \item The method extrapolates the global evolution of structural variables, not the detailed intra-cycle oscillations.
    \item The control function is the central contribution of the paper.
    \item The jump length is automatically shortened when the response becomes nonlinear.
    \item A converged finite element solution after a jump is not necessarily correct.
    \item Accuracy must be checked using state-variable comparison with a reference or controlled error measure.
    \item The method is most suitable for quasi-linear global evolution but can still handle nonlinear regions by reducing or suppressing jumps.
\end{itemize}

\subsection{Limitations}

\begin{itemize}[leftmargin=*]
    \item The method depends on user-defined tolerance values $q_y$.
    \item If $q_y$ is too large, the method may take jumps that are too long and lose accuracy.
    \item If $q_y$ is too small, the method becomes accurate but less efficient.
    \item The approach requires careful selection of control variables.
    \item The method does not automatically guarantee correctness only from finite element convergence.
\end{itemize}

\subsection{Relevance for Thesis Work}

This paper is important because it gives one of the clearest early finite-element cycle-jump algorithms.

For future UMAT or Abaqus cycle-jump implementation, the most useful ideas are:

\begin{itemize}[leftmargin=*]
    \item extract cycle-to-cycle evolution of selected variables,
    \item define a control set,
    \item compute local allowed jump lengths,
    \item take the global minimum jump length,
    \item extrapolate using a Heun-type formula,
    \item reinsert the extrapolated state into Abaqus,
    \item continue the finite element simulation after the jump.
\end{itemize}

\subsection{One-Sentence Summary}

Cojocaru and Karlsson's method accelerates cyclic finite element simulations by computing a few real cycles, estimating the global evolution slope of selected state variables, using a control function to choose a safe jump length, and extrapolating the state with a Heun-type formula before continuing the finite element analysis.

% =========================================================
\section{Nesnas and Saanouni (2000): Cycle-Jumping Scheme for Coupled Damage--Viscoplastic Models}

\subsection{Full Reference}

Nesnas, K. and Saanouni, K. (2000).
\textit{A cycle jumping scheme for numerical integration of coupled damage and viscoplastic models for cyclic loading paths}.
Revue Europ\'eenne des \'El\'ements Finis, 9(8), 865--891.

\subsection{Main Problem Addressed}

This paper develops a cycle-jumping scheme for the numerical integration of coupled damage and viscoplastic constitutive models under cyclic loading.

The problem is that a fully coupled damage--viscoplastic simulation requires the constitutive equations to be integrated through many loading cycles. This is computationally expensive because:

\begin{itemize}[leftmargin=*]
    \item the constitutive equations are nonlinear and stiff,
    \item damage evolves progressively over cycles,
    \item viscoplasticity and damage are coupled,
    \item direct cycle-by-cycle integration may require thousands of cycles.
\end{itemize}

The proposed solution is a \textbf{two time-scale integration scheme}:

\begin{itemize}[leftmargin=*]
    \item \textbf{Small time scale}: implicit integration inside each loading cycle.
    \item \textbf{Large time scale}: explicit cycle-jump integration over many cycles.
\end{itemize}

\subsection{Main Idea of the Cycle-Jump Method}

The fatigue life is divided into three stages:

\begin{enumerate}[leftmargin=*]
    \item \textbf{Hardening stage}: stress increases from cycle to cycle.
    \item \textbf{Stabilized stage}: stress redistribution is slow.
    \item \textbf{Softening stage}: damage grows rapidly and stress decreases until failure.
\end{enumerate}

The stabilized stage is usually the longest part of the fatigue life. Therefore, the method skips some intermediate cycles during this stage.

The central idea is:

\begin{quote}
Integrate accurately inside a few real cycles, estimate how selected variables evolve from cycle to cycle, jump over several cycles, and extrapolate the material state at the new cycle number.
\end{quote}

\subsection{State Variable as a Function of Cycle Number}

Let $T$ be the period of one loading cycle. A particular instant $\tau$ inside each cycle is selected, where

\[
0 \leq \tau \leq T.
\]

The vector of state variables is then written as a function of cycle number $N$:

\[
\bm{y}(N)
=
\bm{y}\big((N-1)T+\tau\big).
\]

Here:

\begin{itemize}[leftmargin=*]
    \item $\bm{y}$ is the vector of state variables,
    \item $N$ is the cycle number,
    \item $T$ is the cycle period,
    \item $\tau$ is the selected point inside the cycle.
\end{itemize}

The problem becomes the determination of a sequence of values

\[
\bm{y}(0 \leq N \leq N_R),
\]

where $N_R$ is the number of cycles to failure.

\subsection{Second-Order Taylor Expansion in Cycle Space}

The state variable after a cycle jump $\Delta N$ is approximated using a Taylor expansion in cycle space:

\[
\bm{y}_{N+\Delta N}
=
\bm{y}_{N}
+
\dot{\bm{y}}_{N}\Delta N
+
\ddot{\bm{y}}_{N}
\frac{\Delta N^2}{2}.
\]

Here:

\begin{itemize}[leftmargin=*]
    \item $\dot{\bm{y}}_N$ is the first pseudo-derivative with respect to cycle number,
    \item $\ddot{\bm{y}}_N$ is the second pseudo-derivative with respect to cycle number,
    \item $\Delta N$ is the number of cycles skipped.
\end{itemize}

This is not a time derivative inside one loading cycle. It is a derivative with respect to the cycle number.

\subsection{First Pseudo-Derivative}

The first pseudo-derivative is approximated as

\[
\dot{\bm{y}}_N
=
\frac{\partial \bm{y}}{\partial N}
\approx
\frac{\bm{y}_{N+1}-\bm{y}_{N}}{1}.
\]

So, in simple form,

\[
\dot{\bm{y}}_N
\approx
\bm{y}_{N+1}-\bm{y}_{N}.
\]

This represents the cycle-to-cycle change of the state variable vector.

\subsection{First-Order Jump Criterion}

The first possibility is to keep only the first-order Taylor term. The jump size is chosen so that the first-order extrapolated change remains small compared with the current variable value:

\[
\left|\dot{\bm{y}}_N\right|\Delta N
\leq
\eta_1
\left|\bm{y}_N\right|.
\]

Solving for the jump size gives

\[
\Delta N
=
\eta_1
\left|
\frac{\bm{y}_N}{\dot{\bm{y}}_N}
\right|.
\]

Here, $\eta_1$ is an accuracy parameter.

\subsubsection*{Meaning}

If the state variable changes slowly from cycle to cycle, then $\left|\dot{\bm{y}}_N\right|$ is small and $\Delta N$ becomes large.

If the state variable changes rapidly, then $\left|\dot{\bm{y}}_N\right|$ is large and $\Delta N$ becomes small.

\subsection{Second-Order Jump Criterion}

The second possibility is to use a second-order Taylor development. The second-order term is controlled relative to the first-order term:

\[
\left|
\ddot{\bm{y}}_N
\right|
\frac{\Delta N^2}{2}
\leq
\eta_2
\left|
\dot{\bm{y}}_N
\right|
\Delta N.
\]

Solving for $\Delta N$ gives

\[
\Delta N
=
2\eta_2
\left|
\frac{\dot{\bm{y}}_N}{\ddot{\bm{y}}_N}
\right|.
\]

Here, $\eta_2$ is another accuracy parameter.

\subsubsection*{Meaning}

This condition controls the curvature of the state-variable evolution.

If the state evolution is almost linear, $\ddot{\bm{y}}_N$ is small and a large jump is allowed.

If the state evolution is strongly nonlinear, $\ddot{\bm{y}}_N$ is large and the jump size is reduced.

\subsection{Global Cycle-Jump Size}

The local jump size is evaluated for the chosen variables at each integration point.

The global cycle-jump increment is then selected as the minimum value:

\[
\Delta N
=
\min_i
\left(
\Delta N_i
\right).
\]

This means the most critical integration point controls the global jump size.

\subsection{Why Not Use All State Variables?}

The classical cycle-jump technique uses all state variables to estimate $\Delta N$.

However, Nesnas and Saanouni point out that this can be difficult because different state variables:

\begin{itemize}[leftmargin=*]
    \item evolve at different rates,
    \item may have very different magnitudes,
    \item may show different degrees of nonlinearity,
    \item may give artificially small or unsuitable jump sizes.
\end{itemize}

Therefore, the paper proposes using only a small number of representative variables that measure the global evolution of the material state.

\subsection{Main Variables Used for Jump Control}

The paper identifies two important variables for controlling the cycle jump:

\[
\Delta L
\qquad \text{and} \qquad
D.
\]

Here:

\begin{itemize}[leftmargin=*]
    \item $\Delta L$ is an important scalar parameter from the asymptotic integration algorithm,
    \item $D$ is the damage variable.
\end{itemize}

The parameter $\Delta L$ is sensitive to the nonlinear viscoplastic evolution, especially during hardening and stabilized stages.

The damage variable $D$ is especially important during the softening stage, where damage grows rapidly near failure.

\subsection{Cycle-Jump Size Based on $\Delta L$ and $D$}

The paper proposes computing two candidate jump sizes:

\[
\Delta N_{\Delta L}
=
\eta
\frac{\Delta L}{\dot{\Delta L}},
\]

\[
\Delta N_D
=
\eta
\frac{D}{\dot{D}}.
\]

The final jump size is

\[
\Delta N
=
\min
\left(
\Delta N_{\Delta L},
\Delta N_D
\right).
\]

\subsubsection*{Interpretation}

The variable $\Delta L$ controls the jump during viscoplastic hardening and stabilization.

The variable $D$ controls the jump during the damage-dominated softening stage.

Taking the minimum makes the algorithm conservative.

\subsection{Overall Cycle-Jump Algorithm}

The cycle-jump algorithm works as follows:

\begin{enumerate}[leftmargin=*]
    \item After each cycle jump, integrate the constitutive equations over a small number of real cycles.
    \item Store the values of the selected jump-control variables.
    \item Compute the first and second pseudo-derivatives in cycle space.
    \item Estimate the cycle-jump increment $\Delta N$.
    \item Extrapolate the state variables over $\Delta N$ cycles.
    \item Continue the full integration for several cycles after the jump.
    \item Repeat until failure.
\end{enumerate}

In the paper, a typical value used for the number of real cycles before jumping is five cycles.

\subsection{Two Versions of the Cycle-Jump Technique}

The paper introduces two versions of the cycle-jump technique.

\subsubsection{CJTV1: Cycle-Jump Technique Version 1}

In the first version, all state variables are extrapolated directly using the Taylor expansion:

\[
\bm{y}_{N+\Delta N}
=
\bm{y}_{N}
+
\dot{\bm{y}}_{N}\Delta N
+
\ddot{\bm{y}}_{N}
\frac{\Delta N^2}{2}.
\]

This version is straightforward but requires storing and extrapolating all state variables.

\subsubsection{CJTV2: Cycle-Jump Technique Version 2}

In the second version, only the selected scalar parameters are extrapolated:

\[
\Delta L,
\qquad
\Delta p,
\qquad
D.
\]

Then, the remaining state variables are reconstructed using the asymptotic integral forms of the constitutive equations.

This reduces the number of variables that must be stored and extrapolated.

\subsection{Small Time-Scale Integration}

Inside each loading cycle, the paper uses an implicit asymptotic integration algorithm.

The coupled damage--viscoplastic model includes:

\begin{itemize}[leftmargin=*]
    \item Cauchy stress $\sigma_{ij}$,
    \item kinematic hardening variables $\alpha^k_{ij}$,
    \item isotropic hardening variable $r$,
    \item damage variable $D$.
\end{itemize}

The recursive update of stress has the form

\[
\begin{aligned}
\sigma_{ij}(t+\Delta t)
&=
X_{ij}(t+\Delta t)
+
\frac{\nu E}{(1+\nu)(1-2\nu)}\varepsilon_{kk}(t)\delta_{ij} \\
&\quad
+
\exp(-\Delta L)
\left[
\sigma_{ij}(t)-X_{ij}(t)
-
\frac{\nu E}{(1+\nu)(1-2\nu)}\varepsilon_{kk}(t)\delta_{ij}
\right]
+
\cdots
\end{aligned}
\]

where $\Delta L$ appears as a central scalar parameter in the integration algorithm.

\subsection{Kinematic Hardening Update}

The kinematic hardening variable is updated as

\[
\alpha^k_{ij}(t+\Delta t)
=
\exp(-\Delta G^k)\alpha^k_{ij}(t)
+
\Delta \varepsilon^p_{ij}
\left[
\frac{1-\exp(-\Delta G^k)}{\Delta G^k}
\right].
\]

The scalar parameter $\Delta G^k$ is

\[
\Delta G^k
=
a_k \sqrt{1-D}\Delta p.
\]

\subsection{Isotropic Hardening Update}

The isotropic hardening variable is updated as

\[
r(t+\Delta t)
=
\exp(-\Delta Q)r(t)
+
\Delta p
\left[
\frac{1-\exp(-\Delta Q)}{\Delta Q}
\right].
\]

The scalar parameter $\Delta Q$ is

\[
\Delta Q
=
b\sqrt{1-D}\Delta p.
\]

\subsection{Damage Update}

The damage variable can be updated using a first-order approximation:

\[
D(t+\Delta t)
=
D(t)
+
\Delta D
=
D(t)
+
\dot{D}(t+\Delta t)\Delta t.
\]

A second-order approximation can also be used:

\[
D(t+\Delta t)
=
D(t)
+
\Delta D
=
D(t)
+
\frac{1}{2}
\left[
\dot{D}(t)+\dot{D}(t+\Delta t)
\right]\Delta t.
\]

The paper shows that the order of damage extrapolation is important, especially near the final damage-dominated stage.

\subsection{Implicit Unknowns in the Small Time-Scale Algorithm}

Although the complete constitutive model contains many variables, the asymptotic algorithm reduces the implicit numerical problem to two main scalar unknowns:

\[
\Delta L
\qquad \text{and} \qquad
\Delta D.
\]

These are solved using a Newton--Raphson procedure.

The nonlinear system is written as

\[
g_1(\Delta L,\Delta D)=0,
\]

\[
g_2(\Delta L,\Delta D)=0.
\]

In compact vector form,

\[
g_i(x_j)=0,
\]

where

\[
x_1=\Delta L,
\qquad
x_2=\Delta D.
\]

The Newton linearization is

\[
\frac{\partial g_i}{\partial x_j}
\delta x_j
=
-g_i(x_j).
\]

The iterative update is

\[
x_j^{k+1}
=
x_j^k
+
\delta x_j^k.
\]

The convergence criteria are

\[
\frac{\|\delta x_j^k\|_2}{\|x_j^k\|_2}
\leq
\varepsilon_1,
\]

\[
\frac{\|g_j^k\|_2}{\|g_j^1\|_2}
\leq
\varepsilon_2.
\]

\subsection{Extrapolation in CJTV2}

In CJTV2, after extrapolating

\[
\Delta L,
\qquad
\Delta p,
\qquad
D,
\]

the remaining state variables are reconstructed.

For example, the stress after the cycle jump is reconstructed using an asymptotic expression of the form

\[
\begin{aligned}
\sigma^{N+\Delta N}_{ij}(t+\Delta t)
&=
\frac{2}{3} C\alpha^{N+\Delta N}_{ij}(t+\Delta t)
+
\frac{\nu E}{(1+\nu)(1-2\nu)}\varepsilon_{kk}(t+\Delta t)\delta_{ij} \\
&\quad
+
\exp(-\Delta L^{N+\Delta N})
\left[
\sigma_{ij}(t)
-
\frac{2}{3}C\alpha_{ij}(t)
-
\frac{\nu E}{(1+\nu)(1-2\nu)}\varepsilon_{kk}(t)\delta_{ij}
\right]
+
\cdots
\end{aligned}
\]

This is why CJTV2 stores fewer variables but still reconstructs the complete material state.

\subsection{Numerical Parameters Used in the Cycle-Jump Method}

The paper introduces several practical parameters:

\begin{itemize}[leftmargin=*]
    \item $\tau$: selected instant inside the cycle,
    \item \texttt{JUMPMIN}: minimum number of cycles computed before the next jump,
    \item $\eta$: accuracy parameter for the jump size,
    \item \texttt{JUMPMAX}: maximum allowed number of jumped cycles.
\end{itemize}

These parameters control the balance between accuracy and speed.

\subsection{Important Numerical Observations}

The paper reports the following key observations:

\begin{itemize}[leftmargin=*]
    \item Large jumps are possible in low-curvature stabilized regions.
    \item Smaller jumps are automatically selected near failure.
    \item Damage requires careful extrapolation because it evolves nonlinearly near the end of life.
    \item Second-order extrapolation of damage gives better fatigue-life prediction than first-order extrapolation.
    \item CJTV2 is efficient because it extrapolates only selected scalar parameters and reconstructs the remaining variables.
\end{itemize}

\subsection{Accuracy and Efficiency}

For proportional loading at Gauss point level, the full calculation gave a reference fatigue life of approximately

\[
N_R = 384 \ \text{cycles}.
\]

Using cycle jumping, the method predicted similar fatigue lives while computing far fewer cycles.

For example, with $\eta=0.1$:

\[
N_R \approx 392 \ \text{cycles}
\]

for CJTV1, and

\[
N_R \approx 391 \ \text{cycles}
\]

for CJTV2.

Thus, both schemes saved about 80\% of the actually computed cycles with a small error in fatigue life.

For a long non-proportional loading example, CJTV2 reduced the computational effort by more than 90\% while maintaining good agreement with the full calculation.

\subsection{Structural-Level Implementation}

At structural level, the cycle jump step is calculated at every Gauss integration point.

The global cycle-jump size is

\[
\Delta N
=
\min_{\text{all Gauss points}}
\left(
\Delta N_i
\right).
\]

This ensures that the most critical material point controls the global jump.

The method was tested on:

\begin{itemize}[leftmargin=*]
    \item a simplified three-bar structure,
    \item a plate with a central circular hole.
\end{itemize}

In both cases, the method reproduced the full cycle-by-cycle solution with significantly reduced computation time.

\subsection{Key Technical Takeaways}

\begin{itemize}[leftmargin=*]
    \item This paper uses a two time-scale strategy.
    \item The small time scale is handled by an implicit asymptotic integration algorithm.
    \item The large time scale is handled by an explicit cycle-jump algorithm.
    \item The cycle-jump variable is not necessarily the full state vector.
    \item The scalar variables $\Delta L$ and $D$ are enough to estimate the jump size effectively.
    \item The global jump size is the minimum jump size over all integration points.
    \item Damage evolution strongly controls the final softening stage.
    \item Second-order extrapolation is important for nonlinear damage evolution.
\end{itemize}

\subsection{Limitations}

\begin{itemize}[leftmargin=*]
    \item The method still depends on user-defined parameters such as $\eta$, \texttt{JUMPMIN}, and \texttt{JUMPMAX}.
    \item The accuracy depends on the selected jump-control variables.
    \item Strongly nonlinear damage evolution near failure requires smaller jumps.
    \item The presented damage model does not include unilateral damage effects, where microcracks may open or close depending on the loading direction.
\end{itemize}

\subsection{Relevance for Thesis Work}

This paper is important because it treats cycle jumping at the constitutive integration level.

Compared with structural cycle-jump methods that extrapolate displacements, stresses, or strains directly, this method focuses on internal material-state evolution.

For future UMAT or constitutive implementation work, the important lessons are:

\begin{itemize}[leftmargin=*]
    \item separate the intra-cycle integration from the cycle-space extrapolation,
    \item use robust implicit integration inside the cycle,
    \item use explicit extrapolation only over cycles,
    \item choose a small set of physically meaningful jump-control variables,
    \item reduce the global jump size near damage localization or failure,
    \item reconstruct the full state after the jump instead of storing every state variable.
\end{itemize}

\subsection{One-Sentence Summary}

Nesnas and Saanouni's method accelerates coupled damage--viscoplastic cyclic simulations by using implicit asymptotic integration inside each cycle and explicit cycle-space extrapolation between cycles, with the jump length controlled mainly by the scalar variables $\Delta L$ and damage $D$.

% =========================================================
\section{Li et al. (2025): Adaptive Cycle-Jump Method for Elasto-Plastic Phase-Field Fatigue Crack Propagation}

\subsection{Full Reference}

Li, J., Hu, Y., Ao, N., Miao, H., Zhang, X., Kang, G., and Kan, Q. (2025).
\textit{An adaptive cycle jump method for elasto-plastic phase field modeling addressing fatigue crack propagation}.
Computer Methods in Applied Mechanics and Engineering, 442, 118074.

\subsection{Main Problem Addressed}

This paper develops an adaptive cycle-jump method for accelerating phase-field simulations of fatigue crack propagation in elasto-plastic materials.

Phase-field fatigue simulations are computationally expensive because:

\begin{itemize}[leftmargin=*]
    \item the crack is represented by a diffused phase-field variable, requiring a fine mesh near the crack path,
    \item fatigue crack propagation normally needs cycle-by-cycle computation,
    \item elasto-plastic crack-tip behavior introduces additional nonlinear constitutive cost,
    \item the crack propagation life is usually not known in advance.
\end{itemize}

The proposed method accelerates the simulation by explicitly extrapolating selected variables over skipped cycles and adaptively choosing the cycle-jump size based on phase-field evolution.

\subsection{Main Idea of the Method}

The key idea is:

\begin{quote}
Compute three complete cycles, use their results to estimate the evolution trend of important phase-field and plasticity variables, jump over $\Delta N$ cycles, and update the variables at cycle $N+\Delta N$.
\end{quote}

Unlike purely elastic phase-field fatigue methods, this method also accounts for plastic strain accumulation at the crack tip.

\subsection{Phase-Field Description of Crack}

The body is described by two primary fields:

\[
\bm{u}
\qquad \text{and} \qquad
\phi.
\]

Here:

\begin{itemize}[leftmargin=*]
    \item $\bm{u}$ is the displacement field,
    \item $\phi$ is the phase-field crack variable,
    \item $\phi = 0$ represents intact material,
    \item $\phi = 1$ represents fully cracked material.
\end{itemize}

The degradation function is

\[
g(\phi)
=
(1-\phi)^2+\kappa,
\]

where $\kappa$ is a small numerical parameter introduced to avoid singularity when the material is fully damaged.

\subsection{Total Potential Energy}

For an elasto-plastic body, the total potential energy is written as

\[
\begin{aligned}
\Psi
&=
\int_{\Omega}
\left[
g(\phi)\psi^e(\bm{\varepsilon}^e)
+
g(\phi)\psi^p(\bm{\varepsilon}^p)
\right] dV \\
&\quad
+
\int_{\Omega}
\frac{G_c}{2l}
\left(
\phi^2+l^2|\nabla \phi|^2
\right)dV
-
\int_{\Omega}
\bm{b}\cdot \bm{u}\,dV
-
\int_{\partial\Omega}
\bm{t}\cdot \bm{u}\,dS.
\end{aligned}
\]

Here:

\begin{itemize}[leftmargin=*]
    \item $\psi^e$ is the elastic strain energy density,
    \item $\psi^p$ is the plastic strain energy density,
    \item $G_c$ is the critical energy release rate,
    \item $l$ is the phase-field length scale,
    \item $\bm{b}$ is the body force,
    \item $\bm{t}$ is the boundary traction.
\end{itemize}

\subsection{Elastic Energy Split}

The tensile and compressive parts of the elastic strain energy density are written as

\[
\psi_e^{\pm}
=
\frac{1}{2}
\lambda
\langle \mathrm{tr}(\bm{\varepsilon}^e)\rangle_{\pm}^2
+
\mu
\mathrm{tr}
\left[
(\bm{\varepsilon}_e^{\pm})^2
\right].
\]

Here:

\begin{itemize}[leftmargin=*]
    \item $\lambda$ and $\mu$ are Lam\'e constants,
    \item $\langle x\rangle_{\pm} = \frac{x\pm |x|}{2}$ is the Macaulay bracket.
\end{itemize}

The energy split is used so that crack growth is mainly driven by tensile deformation, not compression.

\subsection{Plastic Strain Energy Density}

The plastic strain energy density is calculated incrementally as

\[
\psi^p
=
\int_0^T
\bm{\sigma}:\dot{\bm{\varepsilon}}^p\,dt.
\]

This represents the accumulated plastic work density.

\subsection{Fatigue Degradation Function}

Fatigue is introduced by reducing the effective fracture toughness through a fatigue degradation function:

\[
f(\alpha)
=
\begin{cases}
1, & \alpha \leq \alpha_T, \\
\left(
\frac{2\alpha_T}{\alpha+\alpha_T}
\right)^2,
& \alpha > \alpha_T.
\end{cases}
\]

Here:

\begin{itemize}[leftmargin=*]
    \item $\alpha$ is the accumulated fatigue history variable,
    \item $\alpha_T$ is the fatigue threshold.
\end{itemize}

The fatigue threshold is

\[
\alpha_T
=
\frac{G_c}{12l}.
\]

\subsection{Fatigue History Variable}

The fatigue driving quantity is

\[
\bar{\alpha}
=
g(\phi)
\left(
\psi_e^+ + \psi^p
\right).
\]

The accumulated fatigue history is

\[
\alpha
=
\int_0^T
H(\dot{\bar{\alpha}})
|\dot{\bar{\alpha}}|dt.
\]

Here, $H(\cdot)$ is the Heaviside function.

This means fatigue history accumulates only when the fatigue driving quantity increases.

\subsection{Final Energy Functional with Fatigue}

After adding fatigue degradation and elastic energy split, the total energy becomes

\[
\begin{aligned}
\Psi
&=
\int_{\Omega}
\left[
g(\phi)(\psi_e^+ + \psi^p)
+
\psi_e^-
\right]dV \\
&\quad
+
\int_{\Omega}
f(\alpha)
\frac{G_c}{2l}
\left(
\phi^2+l^2|\nabla \phi|^2
\right)dV
-
\int_{\Omega}
\bm{b}\cdot \bm{u}\,dV
-
\int_{\partial\Omega}
\bm{t}\cdot \bm{u}\,dS.
\end{aligned}
\]

\subsection{History Field for Crack Irreversibility}

To enforce crack irreversibility, a history field is introduced:

\[
P
=
\max_{\tau\in[0,t]}
\left\{
\psi_e^+ + \eta \psi^p
\right\}.
\]

Here, $\eta$ controls the contribution of plastic work to the crack-driving history field.

The paper studies cases such as

\[
\eta = 0.1
\qquad \text{and} \qquad
\eta = 1.
\]

\subsection{Governing Equations}

The governing equations are

\[
\begin{cases}
\nabla\cdot \bm{\sigma} + \bm{b} = 0, \\
g'(\phi)P
+
\dfrac{f(\alpha)G_c}{l}\phi
-
G_c l \nabla\cdot
\left[
f(\alpha)\nabla\phi
\right]
=0.
\end{cases}
\]

The paper uses a hybrid phase-field form:

\[
\begin{cases}
\nabla\cdot \bm{\sigma} + \bm{b} = 0, \\
g'(\phi)P
+
\dfrac{f(\alpha)G_c}{l}\phi
-
G_c l \nabla\cdot
\left[
f(\alpha)\nabla\phi
\right]
=0, \\
\phi = 0
\quad \text{if}
\quad
\psi_e^+ < \psi_e^-.
\end{cases}
\]

The hybrid form keeps the displacement equation computationally simpler while still avoiding compression-driven crack growth.

\subsection{Plastic Constitutive Models}

The paper verifies the method using three plasticity models:

\begin{enumerate}[leftmargin=*]
    \item linear isotropic hardening,
    \item power-law hardening,
    \item Armstrong--Frederick cyclic plasticity.
\end{enumerate}

All models are based on the von Mises yield criterion.

\subsection{Yield Function}

The yield function is

\[
F_y
=
\sigma_{\mathrm{eq}} - Q.
\]

Here:

\begin{itemize}[leftmargin=*]
    \item $\sigma_{\mathrm{eq}}$ is the equivalent von Mises stress,
    \item $Q$ is the current yield stress.
\end{itemize}

\subsection{Equivalent Stress for Isotropic Hardening}

For linear and power-law isotropic hardening models,

\[
\sigma_{\mathrm{eq}}
=
\sqrt{
\frac{3}{2}
\bm{s}:\bm{s}
},
\]

where $\bm{s}$ is the deviatoric stress tensor.

\subsection{Plastic Flow Rule}

The plastic strain rate is

\[
\dot{\bm{\varepsilon}}^p
=
\sqrt{\frac{3}{2}}
\dot{\varepsilon}_{\mathrm{eq}}^p
\frac{\bm{s}}{\sqrt{\bm{s}:\bm{s}}}.
\]

The equivalent plastic strain rate is

\[
\dot{\varepsilon}_{\mathrm{eq}}^p
=
\sqrt{
\frac{2}{3}
\dot{\bm{\varepsilon}}^p:\dot{\bm{\varepsilon}}^p
}.
\]

\subsection{Yield Stress for Linear and Power-Law Hardening}

The current yield stress is

\[
Q
=
\begin{cases}
\sigma_y + h\varepsilon_{\mathrm{eq}}^p,
& \text{linear isotropic hardening}, \\
\sigma_y
\left(
1+\dfrac{E}{\sigma_y}\varepsilon_{\mathrm{eq}}^p
\right)^n,
& \text{power-law hardening}.
\end{cases}
\]

Here:

\begin{itemize}[leftmargin=*]
    \item $\sigma_y$ is the initial yield stress,
    \item $h$ is the linear hardening modulus,
    \item $E$ is Young's modulus,
    \item $n$ is the strain hardening exponent.
\end{itemize}

\subsection{Armstrong--Frederick Model}

For the Armstrong--Frederick model, the isotropic hardening rule is

\[
Q
=
\sigma_y
+
(Q_{\mathrm{sat}}-\sigma_y)
\left[
1-\exp(-b\varepsilon_{\mathrm{eq}}^p)
\right].
\]

The equivalent stress is

\[
\sigma_{\mathrm{eq}}
=
\sqrt{
\frac{3}{2}
(\bm{s}-\bm{\theta}):(\bm{s}-\bm{\theta})
}.
\]

The plastic strain rate becomes

\[
\dot{\bm{\varepsilon}}^p
=
\sqrt{\frac{3}{2}}
\dot{\varepsilon}_{\mathrm{eq}}^p
\frac{\bm{s}-\bm{\theta}}
{\sqrt{(\bm{s}-\bm{\theta}):(\bm{s}-\bm{\theta})}}.
\]

The backstress evolution is

\[
\dot{\bm{\theta}}
=
\frac{2}{3}C\dot{\bm{\varepsilon}}^p
-
\gamma \bm{\theta}\dot{\varepsilon}_{\mathrm{eq}}^p.
\]

Here:

\begin{itemize}[leftmargin=*]
    \item $\bm{\theta}$ is the backstress tensor,
    \item $C$ is the kinematic hardening modulus,
    \item $\gamma$ is the dynamic recovery parameter.
\end{itemize}

\subsection{Finite Element Discretization}

The displacement and phase fields are approximated as

\[
\bm{u}
=
\sum_{i=1}^{m}
N_i^u \bm{u}_i,
\qquad
\phi
=
\sum_{i=1}^{m}
N_i^\phi \phi_i.
\]

The strain and phase-field gradient are

\[
\bm{\varepsilon}
=
\sum_{i=1}^{m}
\bm{B}_i^u\bm{u}_i,
\qquad
\nabla \phi
=
\sum_{i=1}^{m}
\bm{B}_i^\phi \phi_i.
\]

\subsection{Residual Equations}

The displacement residual is

\[
\begin{aligned}
\bm{R}_i^u
&=
\int_{\Omega}
(\bm{B}_i^u)^T\bm{\sigma}\,dV
-
\int_{\Omega}
(N_i^u)^T\bm{b}\,dV \\
&\quad
-
\int_{\partial\Omega}
(N_i^u)^T\bm{t}\,dS.
\end{aligned}
\]

The phase-field residual is

\[
\begin{aligned}
R_i^\phi
=
\int_{\Omega}
\left\{
\left[
\frac{f(\alpha)G_c}{l}\phi
+
g'(\phi)P
\right]
N_i^\phi
+
f(\alpha)G_c l
(\bm{B}_i^\phi)^T\nabla\phi
\right\}dV.
\end{aligned}
\]

\subsection{Stiffness Matrices}

The displacement stiffness matrix is

\[
\begin{aligned}
\bm{K}_{ij}^{u}
&=
\frac{\partial \bm{R}_i^u}{\partial \bm{u}_j}
=
\int_{\Omega}
g(\phi)
(\bm{B}_i^u)^T
\bm{D}
\bm{B}_j^u
\,dV.
\end{aligned}
\]

The phase-field stiffness matrix is

\[
\begin{aligned}
K_{ij}^{\phi}
&=
\frac{\partial R_i^\phi}{\partial \phi_j}
=
\int_{\Omega}
\Bigl\{
\left[
\frac{f(\alpha)G_c}{l}
+
g''(\phi)P
\right]
N_i^\phi N_j^\phi \\
&\quad
+
f(\alpha)G_c l
(\bm{B}_i^\phi)^T
\bm{B}_j^\phi
\Bigr\}dV.
\end{aligned}
\]

Here, $\bm{D}$ is the consistent tangent stiffness matrix of the intact elasto-plastic material.

\subsection{Cycle-Jump Variable Vector}

The selected variable vector for cycle jumping is

\[
\bm{\beta}
=
\left[
P,\ \alpha,\ \varepsilon_{\mathrm{eq}}^p
\right].
\]

Here:

\begin{itemize}[leftmargin=*]
    \item $P$ controls the phase-field crack-driving force,
    \item $\alpha$ controls fatigue degradation,
    \item $\varepsilon_{\mathrm{eq}}^p$ accounts for plastic strain accumulation at the crack tip.
\end{itemize}

This is one of the most important choices in the paper.

\subsection{Why These Variables Are Selected}

The method does not extrapolate every field variable. Instead, it updates the quantities that most directly control crack propagation.

\begin{itemize}[leftmargin=*]
    \item $P$ determines the history-dependent crack-driving energy.
    \item $\alpha$ controls the reduction of fracture resistance under fatigue.
    \item $\varepsilon_{\mathrm{eq}}^p$ accounts for plastic accumulation around the crack tip.
\end{itemize}

Thus, the crack propagation path and rate can be predicted without explicitly extrapolating all displacement and phase-field variables.

\subsection{Evolution Rate of $\beta$}

The evolution rate of $\bm{\beta}$ at cycle $N$ is defined as

\[
\bm{v}_N
=
\left.
\frac{\partial \bm{\beta}}{\partial N}
\right|_N.
\]

In discrete form,

\[
\bm{v}_N
=
\bm{\beta}_N-\bm{\beta}_{N-1}.
\]

This is the cycle-to-cycle change of the variable vector.

\subsection{Trapezoidal Estimate of Variable Change}

The change in $\bm{\beta}$ over a cycle jump $\Delta N$ is approximated as

\[
\Delta \bm{\beta}
=
\Delta N
\frac{
\bm{v}_N+\bm{v}_{N+\Delta N}
}{2}.
\]

Therefore,

\[
\bm{\beta}_{N+\Delta N}
=
\bm{\beta}_{N}
+
\Delta N
\frac{
\bm{v}_N+\bm{v}_{N+\Delta N}
}{2}.
\]

\subsection{Earlier Second-Order Taylor Form}

Earlier explicit cycle-jump methods approximate the variable update as

\[
\bm{\beta}_{N+\Delta N}
=
\bm{\beta}_N
+
\Delta N\bm{v}_N
+
\frac{\Delta N^2}{2}\bm{a}_N.
\]

Here,

\[
\bm{a}_N
=
\bm{v}_N-\bm{v}_{N-1}.
\]

This assumes that the evolution rate $\bm{v}$ changes linearly during the jump.

\subsection{Modified Evolution Rate for Nonlinear Evolution}

Li et al. modify the evolution rate to better account for nonlinear evolution:

\[
\bm{v}_{N+\Delta N}
=
\bm{v}_N
+
\Delta N\bm{a}_N
+
\frac{\Delta N^2}{2}
(\bm{a}_N-\bm{a}_{N-1}).
\]

Substituting this into the trapezoidal expression gives

\[
\bm{\beta}_{N+\Delta N}
=
\bm{\beta}_N
+
\Delta N\bm{v}_N
+
\frac{\Delta N^2}{2}\bm{a}_N
+
\frac{\Delta N^3}{4}
(\bm{a}_N-\bm{a}_{N-1}).
\]

\subsection{Expression Using Previous Cycle Values}

Using

\[
\bm{v}_N
=
\bm{\beta}_N-\bm{\beta}_{N-1},
\]

and

\[
\bm{a}_N
=
\bm{v}_N-\bm{v}_{N-1},
\]

the update becomes

\[
\begin{aligned}
\bm{\beta}_{N+\Delta N}
&=
\bm{\beta}_N
+
\Delta N
(\bm{\beta}_N-\bm{\beta}_{N-1}) \\
&\quad
+
\frac{\Delta N^2}{2}
(\bm{\beta}_N-2\bm{\beta}_{N-1}+\bm{\beta}_{N-2}) \\
&\quad
+
\frac{\Delta N^3}{4}
(\bm{\beta}_N-3\bm{\beta}_{N-1}+3\bm{\beta}_{N-2}-\bm{\beta}_{N-3}).
\end{aligned}
\]

\subsection{Adapted Formula Using Three Complete Cycles}

To avoid computing four complete cycles, the paper uses the beginning value of cycle $N-2$:

\[
\begin{aligned}
\bm{\beta}_{q,N+\Delta N}
&=
\bm{\beta}_{q,N}
+
\Delta N
(\bm{\beta}_{q,N}-\bm{\beta}_{q,N-1}) \\
&\quad
+
\frac{\Delta N^2}{2}
(\bm{\beta}_{q,N}-2\bm{\beta}_{q,N-1}+\bm{\beta}_{q,N-2}) \\
&\quad
+
\frac{\Delta N^3}{4}
(\bm{\beta}_{q,N}-3\bm{\beta}_{q,N-1}+3\bm{\beta}_{q,N-2}-\bm{\beta}_{0,N-2}).
\end{aligned}
\]

Here:

\begin{itemize}[leftmargin=*]
    \item subscript $q$ denotes the end of a cycle,
    \item subscript $0$ denotes the beginning of a cycle,
    \item $\bm{\beta}_{q,N}$ is the value at the end of cycle $N$,
    \item $\bm{\beta}_{0,N-2}$ is the value at the beginning of cycle $N-2$.
\end{itemize}

\subsection{Incremental Update Form}

Because the evolution expression is nonlinear, the jump interval is divided into unit intervals. The final update is written as

\[
\bm{\beta}_{q,N+\Delta N}
=
\bm{\beta}_{q,N}
+
\sum_{i=0}^{\Delta N-1}
\Delta \bm{\beta}_i.
\]

The interval increment is

\[
\begin{aligned}
\Delta \bm{\beta}_i
&=
(\bm{\beta}_{q,N+i}-\bm{\beta}_{q,N+i-1}) \\
&\quad
+
\frac{1}{2}
(\bm{\beta}_{q,N+i}-2\bm{\beta}_{q,N+i-1}+\bm{\beta}_{q,N+i-2}) \\
&\quad
+
\frac{1}{4}
(\bm{\beta}_{q,N+i}-3\bm{\beta}_{q,N+i-1}+3\bm{\beta}_{q,N+i-2}-\bm{\beta}_{0,N+i-2}).
\end{aligned}
\]

This recursive update is used during the explicit cycle-jump step.

\subsection{Adaptive Cycle-Jump Size}

The cycle-jump size $\Delta N$ is selected adaptively based on phase-field evolution.

The phase field after a proposed jump is estimated as

\[
\phi_{N+\Delta N}
=
\phi_N
+
\Delta N(\phi_N-\phi_{N-1})
+
\frac{\Delta N^2}{2}
(\phi_N-2\phi_{N-1}+\phi_{N-2}).
\]

The adaptive condition is

\[
\Delta \phi_{N+\Delta N}
\leq
\phi_{\mathrm{cr}}.
\]

Here, $\phi_{\mathrm{cr}}$ is the allowable threshold for the phase-field change.

\subsection{Threshold for Phase-Field Evolution}

The threshold is defined as

\[
\phi_{\mathrm{cr}}
=
\delta_{\mathrm{cr}}\chi_{\mathrm{cr}}\phi.
\]

Here:

\begin{itemize}[leftmargin=*]
    \item $\delta_{\mathrm{cr}}$ is a correction coefficient related to the phase-field evolution rate,
    \item $\chi_{\mathrm{cr}}$ is an upper-limit correction factor,
    \item $\phi$ is the current phase-field value.
\end{itemize}

\subsection{Correction Coefficient}

The correction coefficient is

\[
\delta_{\mathrm{cr},N}
=
\frac{1}{
1+
\left(
\left.
\dfrac{\partial^2\phi}{\partial N^2}
\right|_N
\right)
/
\left(
\left.
\dfrac{\partial \phi}{\partial N}
\right|_N
\right)
}.
\]

In discrete form,

\[
\delta_{\mathrm{cr},N}
=
\frac{\phi_N-\phi_{N-1}}
{
(\phi_N-\phi_{N-1})
+
(\phi_N-2\phi_{N-1}+\phi_{N-2})
}.
\]

\subsection{Final Adaptive Criterion}

Combining the phase-field prediction and threshold gives

\[
\begin{aligned}
&\Delta N(\phi_N-\phi_{N-1})
+
\frac{\Delta N^2}{2}
(\phi_N-2\phi_{N-1}+\phi_{N-2}) \\
&\leq
\frac{\phi_N-\phi_{N-1}}
{
(\phi_N-\phi_{N-1})
+
(\phi_N-2\phi_{N-1}+\phi_{N-2})
}
\chi_{\mathrm{cr}}\phi_N.
\end{aligned}
\]

This criterion is solved iteratively to obtain the allowable cycle-jump size.

\subsection{Element-Wise Jump Size and Global Minimum}

Because the adaptive cycle-jump calculation is performed element-wise in the UEL implementation, each element may produce a different jump size.

The element-wise jump size is stored as

\[
\Delta N(i,e)
=
\Delta N_{i,e},
\]

where $i$ is the node number and $e$ is the element number.

The global jump size is then chosen as

\[
\Delta N_{\min}
=
\min
\left(
\Delta N(i,e)
\right)
-
1.
\]

This makes the most restrictive crack-tip region control the global jump.

\subsection{Complete Adaptive Explicit Cycle-Jump Algorithm}

The algorithm is:

\begin{enumerate}[leftmargin=*]
    \item Compute three complete cycles.
    \item Store
    \[
    \bm{\beta}_{0,N-2},
    \quad
    \bm{\beta}_{q,N-2},
    \quad
    \bm{\beta}_{q,N-1},
    \quad
    \bm{\beta}_{q,N}.
    \]
    \item Compute the phase-field threshold $\phi_{\mathrm{cr}}$.
    \item Initialize trial cycle-jump size $\Delta N=1$.
    \item Estimate $\Delta \phi_{N+\Delta N}$.
    \item Increase $\Delta N$ while
    \[
    \Delta \phi_{N+\Delta N}\leq \phi_{\mathrm{cr}}.
    \]
    \item Stop if the maximum allowed jump size is reached.
    \item Store element-wise jump sizes.
    \item Select the global minimum jump size.
    \item Update $\bm{\beta}$ using the recursive cycle-jump formula.
    \item Solve for the displacement and phase fields after the jump.
    \item Continue with the next acceleration period.
\end{enumerate}

\subsection{Behavior of the Adaptive Jump Size}

The adaptive jump size usually has three stages:

\begin{enumerate}[leftmargin=*]
    \item \textbf{Initial stage}: $\Delta N$ is small because phase-field evolution near the initial crack tip changes rapidly.
    \item \textbf{Stable propagation stage}: $\Delta N$ becomes large because crack growth is relatively stable.
    \item \textbf{Final failure stage}: $\Delta N$ decreases again because crack propagation accelerates rapidly.
\end{enumerate}

Thus, the method automatically balances accuracy and efficiency.

\subsection{Implementation in Abaqus}

The method is implemented in Abaqus using:

\begin{itemize}[leftmargin=*]
    \item \texttt{UEL} for displacement and phase-field equations,
    \item \texttt{UMAT} with CPE4 elements as a post-processing layer,
    \item return mapping for the plastic constitutive models,
    \item staggered solution of displacement and phase field.
\end{itemize}

\subsection{Numerical Validation Cases}

The method is verified using several examples:

\begin{itemize}[leftmargin=*]
    \item single-edge notched tension specimen,
    \item single-edge notched shear specimen,
    \item three-point bending specimen,
    \item three-crack square plate under biaxial loading,
    \item compact tension specimen.
\end{itemize}

\subsection{Main Numerical Results}

The paper shows that the proposed method:

\begin{itemize}[leftmargin=*]
    \item reproduces fatigue crack paths accurately,
    \item predicts fatigue crack propagation life with low relative error,
    \item works for elastic and elasto-plastic materials,
    \item works for different plasticity models,
    \item significantly reduces calculation time,
    \item reproduces Paris-law behavior in the compact tension example.
\end{itemize}

Most reported relative errors are below approximately $5\%$ for the tested cases.

\subsection{Paris Law Verification}

For the compact tension specimen, the paper checks whether the accelerated simulation follows Paris' law:

\[
\frac{da}{dN}
=
L(\Delta K)^m.
\]

The stress intensity factor range is calculated as

\[
\begin{aligned}
\Delta K
&=
\frac{\Delta F}{B\sqrt{W}}
\frac{2+a/W}{(1-a/W)^{1.5}}
\left[
0.886
+
4.64\left(\frac{a}{W}\right)
-
13.32\left(\frac{a}{W}\right)^2
+
14.72\left(\frac{a}{W}\right)^3
-
5.6\left(\frac{a}{W}\right)^4
\right].
\end{aligned}
\]

Here:

\begin{itemize}[leftmargin=*]
    \item $a$ is crack length,
    \item $B$ is specimen thickness,
    \item $W$ is specimen width,
    \item $\Delta F$ is the load range.
\end{itemize}

The paper reports that the accelerated and unaccelerated $\Delta K$--$da/dN$ curves are consistent.

\subsection{Key Technical Takeaways}

\begin{itemize}[leftmargin=*]
    \item The method uses an explicit cycle-jump strategy.
    \item Three complete cycles are computed before each jump.
    \item The jump variables are
    \[
    \bm{\beta}
    =
    [P,\alpha,\varepsilon_{\mathrm{eq}}^p].
    \]
    \item The proposed update formula includes nonlinear evolution of the variable rate.
    \item The adaptive jump size is controlled by phase-field evolution.
    \item The global jump size is the minimum value over elements and nodes.
    \item Crack-tip phase-field evolution naturally controls the jump size.
    \item The method is suitable for elasto-plastic fatigue crack propagation.
\end{itemize}

\subsection{Limitations}

\begin{itemize}[leftmargin=*]
    \item The method has been verified mainly for classical elasto-plastic constitutive models.
    \item More complex viscoplastic models need further study.
    \item Non-proportional multiaxial loading requires further research.
    \item The choice of crack-driving force and plasticity model can affect the result.
    \item The adaptive correction factor $\chi_{\mathrm{cr}}$ influences both accuracy and efficiency.
\end{itemize}

\subsection{Relevance for Thesis Work}

This paper is important because it connects cycle jumping with elasto-plastic phase-field fatigue crack propagation.

For thesis work, the most important lessons are:

\begin{itemize}[leftmargin=*]
    \item cycle-jump methods can be adapted to phase-field fatigue,
    \item the crack-driving variables should be selected carefully,
    \item plastic strain accumulation must be included for elasto-plastic fatigue,
    \item phase-field evolution is a good adaptive criterion for crack-growth simulations,
    \item using only selected variables can reduce computational cost,
    \item adaptive jump control is necessary because crack growth is strongly nonlinear near failure.
\end{itemize}

\subsection{One-Sentence Summary}

Li et al.'s method accelerates elasto-plastic phase-field fatigue crack propagation by explicitly extrapolating the crack-driving variables $P$, $\alpha$, and $\varepsilon_{\mathrm{eq}}^p$ over skipped cycles while adaptively selecting the jump size from the current phase-field evolution.

% =========================================================
\section{Cheng, Hu and Kirka (2022): Cycle-Jump Acceleration Method for Crystal Plasticity High-Cycle Fatigue}

\subsection{Full Reference}

Cheng, J., Hu, X., and Kirka, M. (2022).
\textit{A cycle-jump acceleration method for the crystal plasticity simulation of high cycle fatigue of the metallic microstructure}.
International Journal of Fatigue.

\subsection{Main Problem Addressed}

This paper develops a cycle-jump acceleration method for crystal plasticity finite element simulations of high-cycle fatigue in metallic microstructures.

Direct cycle-by-cycle crystal plasticity finite element simulation is computationally expensive because:

\begin{itemize}[leftmargin=*]
    \item each loading and unloading cycle requires many finite element increments,
    \item each integration point contains many crystal plasticity state variables,
    \item high-cycle fatigue may require thousands or millions of cycles,
    \item microstructure-level stress redistribution evolves slowly over cycles.
\end{itemize}

The proposed method accelerates the simulation by extrapolating the internal crystal plasticity variables over skipped cycles and then using a re-equilibration step to recover the displacement and stress fields.

\subsection{Main Idea of the Method}

The central idea is:

\begin{quote}
Compute a few complete cycles, estimate the evolution rate of the crystal plasticity internal variables, extrapolate these variables over $\Delta N$ skipped cycles, then perform a re-equilibration step so that the finite element displacement and stress fields become consistent with the extrapolated material state.
\end{quote}

Unlike some structural cycle-jump methods, this method does not directly extrapolate displacement or stress. Instead, it extrapolates the internal variables that store the material history.

\subsection{Microstructure-Based Model}

The paper studies an ACMZ aluminium alloy microstructure obtained from EBSD data.

The microstructure contains:

\begin{itemize}[leftmargin=*]
    \item aluminium matrix grains,
    \item Cu-based grain-boundary particles,
    \item grain orientations from EBSD,
    \item finite element mesh representation of the microstructure.
\end{itemize}

The aluminium grains are modeled using crystal plasticity, while the Cu-based particles are treated as elastic.

\subsection{Crystal Plasticity Kinematics}

The total deformation gradient is multiplicatively decomposed into elastic and plastic parts:

\[
\bm{F}
=
\bm{F}^e \bm{F}^p.
\]

Here:

\begin{itemize}[leftmargin=*]
    \item $\bm{F}$ is the total deformation gradient,
    \item $\bm{F}^e$ is the elastic deformation gradient,
    \item $\bm{F}^p$ is the plastic deformation gradient.
\end{itemize}

The plastic velocity gradient is

\[
\bm{L}^p
=
\dot{\bm{F}}^p
(\bm{F}^p)^{-1}.
\]

For crystal plasticity, it is written as the sum of slip-system contributions:

\[
\bm{L}^p
=
\sum_{\alpha}
\dot{\gamma}^{\alpha}
\bm{s}^{\alpha}
\otimes
\bm{m}^{\alpha}.
\]

Here:

\begin{itemize}[leftmargin=*]
    \item $\alpha$ is the slip-system index,
    \item $\dot{\gamma}^{\alpha}$ is the slip rate on slip system $\alpha$,
    \item $\bm{s}^{\alpha}$ is the slip direction,
    \item $\bm{m}^{\alpha}$ is the slip-plane normal.
\end{itemize}

\subsection{Stress Relation}

The second Piola--Kirchhoff stress is obtained from the elastic strain measure:

\[
\bm{S}
=
\mathbb{C}^e
:
\bm{E}^e.
\]

Here:

\begin{itemize}[leftmargin=*]
    \item $\bm{S}$ is the second Piola--Kirchhoff stress,
    \item $\mathbb{C}^e$ is the anisotropic elastic stiffness tensor,
    \item $\bm{E}^e$ is the elastic Green strain.
\end{itemize}

The Cauchy stress is obtained by pushing forward the stress measure to the current configuration.

\subsection{Power-Law Slip Rule}

The slip rate on slip system $\alpha$ is governed by a power-law relation:

\[
\dot{\gamma}^{\alpha}
=
\dot{\gamma}_0
\left|
\frac{
\tau^{\alpha}-\chi^{\alpha}
}{
g^{\alpha}
}
\right|^{1/m}
\mathrm{sign}
\left(
\tau^{\alpha}-\chi^{\alpha}
\right).
\]

Here:

\begin{itemize}[leftmargin=*]
    \item $\dot{\gamma}_0$ is the reference slip rate,
    \item $\tau^{\alpha}$ is the resolved shear stress,
    \item $\chi^{\alpha}$ is the slip-system backstress,
    \item $g^{\alpha}$ is the slip resistance,
    \item $m$ is the strain-rate sensitivity parameter.
\end{itemize}

\subsection{Accumulated Slip}

The accumulated slip is

\[
\gamma
=
\int_0^t
\sum_{\alpha}
\left|
\dot{\gamma}^{\alpha}
\right|
\,d\xi.
\]

This variable measures the total accumulated slip activity in the crystal.

\subsection{Slip Resistance Hardening}

The slip resistance evolves according to a saturation-type hardening law:

\[
\begin{aligned}
g^{\alpha}
&=
g_0^{\alpha}
+
\left(
g_s^{\alpha}
-
g_0^{\alpha}
\right)
\left[
1
-
\exp
\left(
-
\frac{
h_0^{\alpha}\gamma
}{
g_s^{\alpha}
}
\right)
\right].
\end{aligned}
\]

Here:

\begin{itemize}[leftmargin=*]
    \item $g_0^{\alpha}$ is the initial slip resistance,
    \item $g_s^{\alpha}$ is the saturation slip resistance,
    \item $h_0^{\alpha}$ is the hardening parameter,
    \item $\gamma$ is the accumulated slip.
\end{itemize}

\subsection{Hall--Petch Grain-Size Dependence}

The initial slip resistance includes grain-size dependence through a Hall--Petch type relation:

\[
g_0^{\alpha}
=
\bar{g}_0^{\alpha}
+
\frac{k}{\sqrt{d^{\alpha}}}.
\]

Here:

\begin{itemize}[leftmargin=*]
    \item $\bar{g}_0^{\alpha}$ is the base initial slip resistance,
    \item $k$ is the Hall--Petch coefficient,
    \item $d^{\alpha}$ is the equivalent grain size.
\end{itemize}

\subsection{Armstrong--Frederick Backstress Evolution}

The slip-system backstress evolves according to an Armstrong--Frederick type law:

\[
\dot{\chi}^{\alpha}
=
c\dot{\gamma}^{\alpha}
-
d\chi^{\alpha}
\left|
\dot{\gamma}^{\alpha}
\right|.
\]

Here:

\begin{itemize}[leftmargin=*]
    \item $c$ controls kinematic hardening,
    \item $d$ controls dynamic recovery,
    \item $\chi^{\alpha}$ shifts the slip response and stores cyclic hardening memory.
\end{itemize}

\subsection{Internal Variable Set for Cycle Jumping}

The crystal plasticity material state is described by the internal variable set

\[
\bm{\mathcal{S}}
=
\left\{
\bm{F}^p,\ \gamma,\ \chi^{\alpha}
\right\}.
\]

Here:

\begin{itemize}[leftmargin=*]
    \item $\bm{F}^p$ stores permanent plastic deformation,
    \item $\gamma$ stores accumulated slip,
    \item $\chi^{\alpha}$ stores the slip-system backstress.
\end{itemize}

This internal variable set is responsible for the path dependence of the crystal plasticity finite element solution.

\subsection{Core Principle of the Cycle-Jump Method}

The cycle-jump procedure consists of five main steps:

\begin{enumerate}[leftmargin=*]
    \item Compute a sufficient number of complete cycles.
    \item Evaluate the evolution rate of the internal variables at each integration point.
    \item Extrapolate the internal variables over $\Delta N$ skipped cycles.
    \item Replace the internal variables in the UMAT by the extrapolated values.
    \item Perform a re-equilibration step to recover consistent displacement and stress fields.
\end{enumerate}

\subsection{Two Key Postulations}

The method relies on two important assumptions.

\subsubsection*{Postulation 1}

The internal variables evolve slowly and smoothly enough during high-cycle plasticity that they can be extrapolated over a limited number of cycles without significant loss of accuracy.

\subsubsection*{Postulation 2}

If the extrapolated internal variables are accurate, then the displacement and stress fields after the skipped cycles can be recovered accurately by a finite element re-equilibration step.

\subsection{Why Re-Equilibration Is Needed}

During the skipped cycles, plastic deformation accumulates. The extrapolated internal variables represent this accumulated plastic state, but the total displacement field has not yet been updated.

Therefore, after the jump:

\begin{itemize}[leftmargin=*]
    \item the plastic strain state is updated,
    \item the total deformation field is still inconsistent,
    \item elastic strain and stress are not yet equilibrated.
\end{itemize}

The re-equilibration step fixes the extrapolated internal variables and lets the implicit finite element solver find a compatible displacement and stress field.

This is why the method is naturally suited to Abaqus/Standard and implicit finite element simulation.

\subsection{Taylor Expansion of Internal Variables}

At integration point $i$, the internal variable array after a cycle jump is approximated by a Taylor expansion in cycle space:

\[
\bm{\mathcal{S}}_{i}^{N+\Delta N}
=
\bm{\mathcal{S}}_{i}^{N}
+
\Delta N
\left(
\bm{\mathcal{S}}_{i}
\right)'
+
\frac{\Delta N^2}{2}
\left(
\bm{\mathcal{S}}_{i}
\right)''
+
O(\Delta N^3).
\]

Here:

\begin{itemize}[leftmargin=*]
    \item $N$ is the current cycle number,
    \item $\Delta N$ is the number of skipped cycles,
    \item $\bm{\mathcal{S}}_{i}^{N}$ is the internal state at cycle $N$,
    \item $(\bm{\mathcal{S}}_i)'$ is the first derivative with respect to cycle number,
    \item $(\bm{\mathcal{S}}_i)''$ is the second derivative with respect to cycle number.
\end{itemize}

\subsection{Backward Difference Approximation}

The first derivative is approximated by a backward difference:

\[
\left(
\bm{\mathcal{S}}_i
\right)'
\approx
\bm{\mathcal{S}}_i^N
-
\bm{\mathcal{S}}_i^{N-1}.
\]

The second derivative is approximated as

\[
\left(
\bm{\mathcal{S}}_i
\right)''
\approx
\bm{\mathcal{S}}_i^N
-
2\bm{\mathcal{S}}_i^{N-1}
+
\bm{\mathcal{S}}_i^{N-2}.
\]

Therefore, at least three complete cycles are required before each jump to evaluate the first and second cycle-space derivatives.

\subsection{Linear Extrapolation Used in the Paper}

Although the Taylor expansion contains higher-order terms, the implemented method uses a backward Eulerian linear approximation:

\[
\bm{\mathcal{S}}_i^{N+\Delta N}
=
\bm{\mathcal{S}}_i^N
+
\Delta N
\left(
\bm{\mathcal{S}}_i
\right)'.
\]

This is accurate only if the neglected second-order term is sufficiently small.

\subsection{Adaptive Error Criterion}

To control the jump length, the second-order term is required to be small compared with the current value of the internal variable:

\[
\frac{\Delta N^2}{2}
\left|
\left(
\bm{\mathcal{S}}_i
\right)''
\right|
\leq
\left|
\bm{\mathcal{S}}_i^N
\right|
\epsilon.
\]

Here, $\epsilon$ is the accuracy control parameter.

\subsection{Allowed Local Jump Size}

Solving the above inequality gives the allowable local jump size:

\[
\Delta N_i
\leq
\sqrt{
\frac{
2
\left|
\bm{\mathcal{S}}_i^N
\right|
\epsilon
}{
\left|
\left(
\bm{\mathcal{S}}_i
\right)''
\right|
}
}.
\]

This equation is the main adaptive jump-size criterion in the paper.

\subsubsection*{Interpretation}

If the second derivative is small,

\[
\left|
\left(
\bm{\mathcal{S}}_i
\right)''
\right|
\rightarrow 0,
\]

then the state variable evolves almost linearly and a large jump is allowed.

If the second derivative is large, the material state is changing nonlinearly and the jump size must be reduced.

\subsection{Global Cycle-Jump Size}

The local allowable jump size is calculated at every integration point.

The global jump size is selected as

\[
\Delta N
=
\min_i
\left\lfloor
\Delta N_i
\right\rfloor.
\]

An upper bound is also imposed:

\[
\Delta N
\leq
\Delta N_{\max}.
\]

Therefore, the most rapidly evolving integration point controls the global cycle jump.

\subsection{Complete Cycle-Jump Algorithm}

The full algorithm is:

\begin{enumerate}[leftmargin=*]
    \item Start at cycle $N$.
    \item Compute at least three full cycles explicitly.
    \item Store the internal variables
    \[
    \bm{\mathcal{S}}^N,
    \quad
    \bm{\mathcal{S}}^{N-1},
    \quad
    \bm{\mathcal{S}}^{N-2}.
    \]
    \item Compute
    \[
    (\bm{\mathcal{S}})'
    \quad \text{and} \quad
    (\bm{\mathcal{S}})''.
    \]
    \item Calculate the local allowable jump size $\Delta N_i$ at each integration point.
    \item Select the global value
    \[
    \Delta N
    =
    \min_i
    \left\lfloor
    \Delta N_i
    \right\rfloor.
    \]
    \item Limit the jump by
    \[
    \Delta N \leq \Delta N_{\max}.
    \]
    \item Extrapolate the internal variables:
    \[
    \bm{\mathcal{S}}^{N+\Delta N}
    =
    \bm{\mathcal{S}}^N
    +
    \Delta N(\bm{\mathcal{S}})'.
    \]
    \item Replace the UMAT state variables by the extrapolated values.
    \item Run the re-equilibration virtual increment.
    \item Update the effective cycle number:
    \[
    N \leftarrow N+\Delta N.
    \]
    \item Continue with the next block of cycles.
\end{enumerate}

\subsection{Known-State-Variable Cycle-Jump Validation}

The paper first validates the method using a known-state-variable cycle-jump test.

In this test:

\begin{itemize}[leftmargin=*]
    \item a normal cycle-by-cycle reference simulation is first performed,
    \item internal variables are stored at every cycle,
    \item a second simulation skips cycles by reading the already-known reference internal variables,
    \item the result is compared with the reference solution.
\end{itemize}

This test checks whether accurate internal variables are sufficient to recover the correct displacement and stress fields through re-equilibration.

The result shows that if the internal variables are known accurately, the skipped-cycle displacement and stress fields can be recovered accurately.

\subsection{Implementation in Abaqus/Standard}

The method is implemented in Abaqus/Standard using user subroutines:

\begin{itemize}[leftmargin=*]
    \item \texttt{UEXTERNALDB},
    \item \texttt{UMAT},
    \item \texttt{DLOAD},
    \item \texttt{DISP}.
\end{itemize}

\subsection{Role of \texttt{UEXTERNALDB}}

The \texttt{UEXTERNALDB} subroutine is used to manage global internal-variable arrays.

Its main tasks are:

\begin{itemize}[leftmargin=*]
    \item initialize arrays for internal variables at all integration points,
    \item store internal variables from previous cycles,
    \item compute first and second cycle-space derivatives,
    \item determine the jump size $\Delta N$,
    \item extrapolate internal variables after the jump,
    \item update the effective cycle number and effective time.
\end{itemize}

\subsection{Role of \texttt{UMAT}}

The \texttt{UMAT} subroutine performs the crystal plasticity constitutive calculation.

During ordinary cycles:

\[
\texttt{UMAT}
\quad
\rightarrow
\quad
\text{regular crystal plasticity update}.
\]

After a cycle jump:

\[
\texttt{UMAT}
\quad
\rightarrow
\quad
\text{read extrapolated internal variables}.
\]

Then the re-equilibration step is performed by fixing the extrapolated state variables and computing stress and tangent stiffness for virtual increments.

\subsection{Role of \texttt{DLOAD} and \texttt{DISP}}

For load-controlled cyclic plasticity:

\begin{itemize}[leftmargin=*]
    \item \texttt{DLOAD} applies periodic loading based on effective time,
    \item during re-equilibration, zero load increment is applied.
\end{itemize}

For displacement-controlled cyclic plasticity:

\begin{itemize}[leftmargin=*]
    \item \texttt{DISP} applies periodic displacement based on effective time,
    \item during re-equilibration, zero displacement increment is applied.
\end{itemize}

\subsection{Effective Cycle Number and Virtual Time}

The paper distinguishes between:

\begin{itemize}[leftmargin=*]
    \item Abaqus step time,
    \item effective cycle number,
    \item effective simulation time.
\end{itemize}

After a jump,

\[
N_{\mathrm{eff}}
\leftarrow
N_{\mathrm{eff}}+\Delta N.
\]

The re-equilibration increment is called a virtual increment because it is not a physical loading increment. It is only used to restore equilibrium after updating the internal material state.

\subsection{Acceleration Behavior}

The jump size is small during early cycles because the material state evolves rapidly during initial plastic redistribution.

As the material response stabilizes, the internal variables evolve more smoothly and the allowable jump size increases.

Therefore:

\[
\text{rapid evolution}
\quad
\Rightarrow
\quad
\Delta N \text{ small},
\]

\[
\text{stable evolution}
\quad
\Rightarrow
\quad
\Delta N \text{ large}.
\]

\subsection{Validation Cases}

The paper validates the method under several conditions:

\begin{itemize}[leftmargin=*]
    \item load-controlled cyclic plasticity with peak stress of 120 MPa and $R=0$,
    \item higher load magnitude with peak stress of 150 MPa,
    \item negative load ratio case with $R=-0.8$,
    \item strain-controlled cyclic plasticity,
    \item elastic shakedown condition,
    \item dislocation-density-based crystal plasticity with SSD and GND.
\end{itemize}

\subsection{Important Numerical Findings}

The paper reports the following trends:

\begin{itemize}[leftmargin=*]
    \item At lower load, the acceleration efficiency is very high because the internal variables evolve smoothly after initial stabilization.
    \item At higher load, the stress and internal variables evolve faster, so the acceleration efficiency is lower.
    \item A negative $R$ ratio causes stronger nonlinear kinematic hardening and reduces the acceleration efficiency.
    \item Strain-controlled cyclic plasticity can also be accelerated successfully.
    \item Elastic shakedown reduces acceleration efficiency during the transition stage, but the acceleration rate recovers after the microstructure becomes mostly elastic.
    \item Dislocation-density-based crystal plasticity can also be accelerated, although GND effects may reduce efficiency because they create localized deformation near grain boundaries.
\end{itemize}

\subsection{Accuracy Measures}

The paper assesses accuracy by comparing:

\begin{itemize}[leftmargin=*]
    \item displacement evolution,
    \item von Mises stress distribution,
    \item accumulated slip evolution,
    \item mean absolute error in the stress field.
\end{itemize}

For the main validation cases, the accelerated results show good agreement with the cycle-by-cycle reference solution.

\subsection{Key Technical Takeaways}

\begin{itemize}[leftmargin=*]
    \item The method accelerates CPFE high-cycle fatigue simulations by extrapolating internal variables, not external response quantities.
    \item The internal variables used are
    \[
    \bm{\mathcal{S}}
    =
    \{\bm{F}^p,\gamma,\chi^\alpha\}.
    \]
    \item A backward Eulerian linear extrapolation is used.
    \item The jump size is controlled by the second cycle-space derivative of the internal variables.
    \item The global jump size is the minimum over all integration points.
    \item Re-equilibration is essential after each jump.
    \item The method is well suited for Abaqus/Standard because re-equilibration requires an implicit solver.
\end{itemize}

\subsection{Limitations}

\begin{itemize}[leftmargin=*]
    \item The method assumes that internal variables evolve smoothly enough over skipped cycles.
    \item Strongly nonlinear evolution reduces the allowable jump size.
    \item Reversing an inaccurate jump is difficult in commercial Abaqus.
    \item The method is currently tied to implicit finite element simulation.
    \item More complex dislocation-density or higher-order GND models may reduce acceleration efficiency.
\end{itemize}

\subsection{Relevance for Thesis Work}

This paper is highly relevant for cycle-jump implementation in user material subroutines because it clearly shows how cycle jumping can be performed at the internal-variable level.

For thesis work, the most important lessons are:

\begin{itemize}[leftmargin=*]
    \item do not extrapolate only stress or displacement;
    \item extrapolate the material history variables stored inside the constitutive model;
    \item use previous cycles to compute cycle-space derivatives;
    \item choose the jump size based on an error criterion;
    \item after jumping, perform re-equilibration to recover a consistent stress and displacement field;
    \item use global arrays or external databases to manage state variables across cycles.
\end{itemize}

\subsection{One-Sentence Summary}

Cheng, Hu and Kirka's method accelerates high-cycle crystal plasticity fatigue simulations by extrapolating the internal variables $\bm{F}^p$, accumulated slip $\gamma$, and slip-system backstress $\chi^\alpha$ over skipped cycles, then performing an implicit re-equilibration step in Abaqus/Standard to recover the consistent displacement and stress fields.

% =========================================================
\section{Comparative Summary of Cycle-Jump Techniques}

\begin{landscape}
\begin{table}[htbp]
\centering
\scriptsize
\renewcommand{\arraystretch}{1.35}
\begin{tabularx}{\linewidth}{p{3.1cm} Y Y Y Y}
\toprule
\textbf{Comparison Point}
&
\textbf{Cojocaru and Karlsson (2006)}
&
\textbf{Nesnas and Saanouni (2000)}
&
\textbf{Li et al. (2025)}
&
\textbf{Cheng, Hu and Kirka (2022)}
\\
\midrule

\textbf{Main objective}
&
Accelerate finite element simulations of cyclically loaded structures with slowly evolving structural or material properties.
&
Accelerate numerical integration of coupled damage--viscoplastic constitutive equations under cyclic loading.
&
Accelerate elasto-plastic phase-field fatigue crack propagation simulations.
&
Accelerate crystal plasticity finite element simulations for high-cycle fatigue of metallic microstructures.
\\

\midrule

\textbf{Level of application}
&
Structural finite element level.
&
Constitutive/material-point level and structural finite element level.
&
Phase-field finite element level using UEL/UMAT implementation.
&
Crystal plasticity finite element level using UMAT and global state-variable storage.
\\

\midrule

\textbf{Main physical problem}
&
Cyclic structural evolution, especially thermal barrier coating evolution, transformation strain, oxidation-growth type problems, and slowly changing material properties.
&
Low-cycle fatigue with coupled viscoplasticity and damage evolution.
&
Fatigue crack propagation in elastic and elasto-plastic materials using phase-field fracture.
&
High-cycle fatigue in metallic microstructures where grain orientation, slip, particles, and local stress redistribution are important.
\\

\midrule

\textbf{Variables used for cycle jump}
&
General finite element variables such as displacement, stress, strain, and possibly equivalent plastic strain. In the examples, displacement components $u_1$ and $u_2$ are mainly used as control variables.
&
Representative internal variables from the asymptotic constitutive integration, mainly $\Delta L$ and damage $D$.
&
Selected phase-field and plasticity variables:
\[
\bm{\beta}=[P,\alpha,\varepsilon_{\mathrm{eq}}^p].
\]
Here, $P$ is the crack-driving history field, $\alpha$ is the fatigue history variable, and $\varepsilon_{\mathrm{eq}}^p$ is equivalent plastic strain.
&
Crystal plasticity internal variables:
\[
\bm{\mathcal{S}}=\{\bm{F}^p,\gamma,\chi^\alpha\}.
\]
Here, $\bm{F}^p$ is plastic deformation gradient, $\gamma$ is accumulated slip, and $\chi^\alpha$ is slip-system backstress.
\\

\midrule

\textbf{Type of extrapolation}
&
Heun-type extrapolation using the current and predicted evolution slopes.
&
Taylor expansion in cycle space. Two versions are discussed: direct extrapolation of all state variables or extrapolation of selected scalar variables followed by reconstruction.
&
Explicit nonlinear cycle-jump extrapolation. The method modifies the evolution-rate prediction to include nonlinear evolution of $\bm{\beta}$.
&
Backward Eulerian linear extrapolation of crystal plasticity internal variables.
\\

\midrule

\textbf{Basic evolution-rate idea}
&
Uses slopes between characteristic values from consecutive cycles:
\[
s_{12}=\frac{y_1-y_2}{\Delta t_{\mathrm{cycle}}},
\qquad
s_{23}=\frac{y_2-y_3}{\Delta t_{\mathrm{cycle}}}.
\]
&
Uses pseudo-derivatives of state variables with respect to cycle number:
\[
\dot{\bm{y}}_N \approx \bm{y}_{N+1}-\bm{y}_N.
\]
&
Uses cycle-to-cycle evolution rate:
\[
\bm{v}_N=\bm{\beta}_N-\bm{\beta}_{N-1}.
\]
Also uses acceleration of this evolution rate.
&
Uses first and second backward differences of internal variables:
\[
(\bm{\mathcal{S}})' \approx \bm{\mathcal{S}}^N-\bm{\mathcal{S}}^{N-1},
\]
\[
(\bm{\mathcal{S}})'' \approx \bm{\mathcal{S}}^N-2\bm{\mathcal{S}}^{N-1}+\bm{\mathcal{S}}^{N-2}.
\]
\\

\midrule

\textbf{Jump-size control}
&
Jump size is controlled by limiting the change in predicted slope:
\[
\Delta t^M_{y,\mathrm{jump}}
=
q_y\Delta t_{\mathrm{cycle}}
\frac{|s_{12}|}{|s_{12}-s_{23}|}.
\]
&
Jump size is controlled mainly by $\Delta L$ and $D$:
\[
\Delta N=\min(\Delta N_{\Delta L},\Delta N_D).
\]
&
Jump size is adaptively controlled by phase-field evolution:
\[
\Delta \phi_{N+\Delta N}\leq \phi_{\mathrm{cr}}.
\]
The most restrictive element controls the global jump.
&
Jump size is controlled by requiring the second-order term to remain small:
\[
\frac{\Delta N^2}{2}|(\bm{\mathcal{S}}_i)''|
\leq
|\bm{\mathcal{S}}_i^N|\epsilon.
\]
\\

\midrule

\textbf{Global jump selection}
&
Minimum admissible jump over all selected variables and material points/nodes.
&
Minimum jump over selected control variables and integration points.
&
Minimum jump over all elements/nodes:
\[
\Delta N_{\min}=\min(\Delta N(i,e))-1.
\]
&
Minimum admissible jump over all integration points:
\[
\Delta N=\min_i \lfloor \Delta N_i \rfloor.
\]
\\

\midrule

\textbf{Need for complete cycles before jump}
&
Needs at least a few complete cycles to estimate the global evolution slope.
&
Computes several real cycles after each jump, for example five cycles, to stabilize the solution and estimate pseudo-derivatives.
&
Uses three complete cycles before each acceleration period.
&
Needs at least three complete cycles to compute first and second cycle-space derivatives.
\\

\midrule

\textbf{How accuracy is maintained}
&
A control function limits slope changes. If the response becomes nonlinear, the jump length automatically becomes small.
&
Accuracy is maintained by controlling $\Delta L$ and $D$, and by computing real cycles between jumps.
&
Accuracy is maintained by adaptive phase-field thresholding. The jump size decreases near rapid crack growth or final failure.
&
Accuracy is maintained by checking the second derivative of internal variables. Strong nonlinear evolution reduces the jump size.
\\

\midrule

\textbf{Treatment after jump}
&
The extrapolated state is imposed into the next finite element analysis using subroutines such as UMAT, DISP, SDVINI, and UEXPAN.
&
The extrapolated variables are used to update or reconstruct the full constitutive state.
&
The variables $P$, $\alpha$, and $\varepsilon_{\mathrm{eq}}^p$ are updated, then the displacement and phase-field equations are solved.
&
The extrapolated internal variables are inserted into UMAT, followed by a re-equilibration virtual increment in Abaqus/Standard.
\\

\midrule

\textbf{Most suitable material/model type}
&
General structural FE models with slowly evolving material properties or boundary conditions.
&
Coupled damage--viscoplastic constitutive models, especially low-cycle fatigue and damage accumulation.
&
Elasto-plastic phase-field fatigue fracture models, especially crack-growth simulations.
&
Crystal plasticity models for microstructure-sensitive high-cycle fatigue.
\\

\midrule

\textbf{Best suited application}
&
Thermal barrier coatings, oxidation growth, transformation strain, cyclic structural evolution, and quasi-linear cyclic degradation.
&
UMAT-type damage--viscoplastic fatigue simulations, material-point fatigue integration, and fatigue life prediction with damage.
&
Abaqus UEL/UMAT phase-field fatigue crack propagation, SENT, SENS, bending, CT, and complex crack-path problems.
&
CPFE simulation of metallic microstructures, grain-level cyclic plasticity, slip accumulation, stress redistribution, SSD/GND-based fatigue models.
\\

\midrule

\textbf{Main advantage}
&
Simple and general finite element cycle-jump framework with automatic jump-length control.
&
Efficient two-time-scale integration and reduction of stored variables by using representative variables.
&
Adaptive method directly connected to phase-field crack evolution; suitable for elastic and elasto-plastic crack propagation.
&
Very efficient for high-cycle CPFE because it extrapolates only internal variables and uses re-equilibration to recover stress/displacement.
\\

\midrule

\textbf{Main limitation}
&
Finite element convergence after a jump does not guarantee physical accuracy; careful control variables are needed.
&
Accuracy depends strongly on selected variables such as $\Delta L$ and $D$, and on damage extrapolation near failure.
&
Adaptive correction factor $\chi_{\mathrm{cr}}$ affects accuracy and efficiency; more complex viscoplastic and multiaxial cases still need further work.
&
Strongly nonlinear internal-variable evolution reduces acceleration; reversing a bad jump is difficult in commercial Abaqus.
\\

\bottomrule
\end{tabularx}
\caption{Comparison of the four cycle-jump techniques based on their level of application, extrapolated variables, jump-control strategy, suitable model type, advantages, and limitations.}
\label{tab:cycle_jump_comparison_papers_as_columns}
\end{table}
\end{landscape}

% =========================================================
\section{Next Paper}

% Add the next paper summary here in the same file.

% =========================================================
\begin{thebibliography}{9}

\bibitem{cojocaru2006}
D. Cojocaru and A. M. Karlsson,
``A simple numerical method of cycle jumps for cyclically loaded structures,''
\textit{International Journal of Fatigue},
vol. 28, pp. 1677--1689, 2006.

\bibitem{cheng2022}
J. Cheng, X. Hu, and M. Kirka,
``A cycle-jump acceleration method for the crystal plasticity simulation of high cycle fatigue of the metallic microstructure,''
\textit{International Journal of Fatigue}, 2022.

\end{thebibliography}

\end{document}



